\subsection{An exact optimizer comparison from information geometry}
\label{app:theory-natural-gradient}

Changing the optimization metric lets us isolate correctness learning from
changes among correct alternatives. The next result shows both effects in
separate equations for $P$ and $q$.


Natural gradient on the categorical simplex uses the Fisher metric, also
called the Shahshahani metric in evolutionary dynamics
\citep[Sections~2.3 and~3.1]{harper2009informationgeometry}. Its gradient
formula gives an exact comparison with the Euclidean-logit collapse result
in Theorem~\ref{thm:grpo-collapse}. We derive the specialization below so
that the dependence on optimizer geometry is explicit.

\begin{theorem}[Natural-gradient selection and replay]
\label{ext:thm-natural-replay}
Retain the finite categorical setup and start from an interior distribution
$p(0)$ with $0<P(0)<1$. Let $\Psi\in C^2([0,1])$ satisfy
$c(P):=\Psi'(P)>0$ on $[0,1]$. Choose a nonempty fixed bank
$\mathcal B\subseteq\mathcal C$, positive weights $w_b$ summing to one over
$\mathcal B$, and extend them by zero to a vector $\bar w$ on
$\mathcal A$. Exact Fisher natural-gradient ascent of
\[
 \mathcal J(p)=\Psi(P)+\rho\sum_{b\in\mathcal B}w_b\log p_b,
 \qquad \rho\ge0,
\]
induces
\begin{align}
 \dot p_i&=c(P)p_i(\mathbf 1\{i\in\mathcal C\}-P)
                     +\rho(\bar w_i-p_i),
 \label{ext:eq-natural-p}\\
 \dot P&=(1-P)(c(P)P+\rho),
 &\dot q_c&=\frac{\rho}{P}(\bar w_c-q_c).
 \label{ext:eq-natural-Pq}
\end{align}
For $\rho=0$, $P(t)\to1$ while $q(t)=q(0)$: correctness learning preserves
all relative probabilities among correct modes. For $\rho>0$, define
$A(t)=\int_0^t\rho/P(s)\,ds$. Every correct mode satisfies
\begin{equation}
 q_c(t)=\bar w_c+(q_c(0)-\bar w_c)e^{-A(t)},
 \qquad A(t)\ge\rho t.
 \label{ext:eq-natural-conditional-solution}
\end{equation}
Furthermore,
\begin{equation}
 1-P(t)\le(1-P(0))e^{-\rho t},
 \qquad
 p_b(t)\ge P(0)\min\{q_b(0),w_b\}>0
 \quad (b\in\mathcal B,t\ge0).
 \label{ext:eq-natural-retention}
\end{equation}
Thus full-support replay sends the conditional distribution to its bank
target, whereas an incomplete bank still loses unbanked correct modes in the
limit.
\end{theorem}

\begin{proof}
For tangent vectors $u,v$ with zero coordinate sum, the categorical Fisher
metric is $g_p(u,v)=\sum_i u_iv_i/p_i$. The tangent vector
\[
 v_i=p_i\left(\partial_i\mathcal J-
                  \sum_jp_j\partial_j\mathcal J\right)
\]
has zero sum and satisfies $g_p(v,u)=\sum_i\partial_i\mathcal J\,u_i$
for every tangent $u$. It is therefore the natural gradient. Here
\[
 \partial_i\mathcal J=c(P)\mathbf 1\{i\in\mathcal C\}
                   +\rho\bar w_i/p_i,
 \qquad
 \sum_i p_i\partial_i\mathcal J=c(P)P+\rho,
\]
which proves Equation~\eqref{ext:eq-natural-p}. Summing over correct
categories gives the $P$ equation. Differentiating $q_c=p_c/P$ then gives
the $q$ equation in~\eqref{ext:eq-natural-Pq}; the correctness terms cancel.

The $P$ equation makes $P$ nondecreasing. If $\rho=0$, $q$ is constant;
if $\rho>0$, solving its linear equation gives
Equation~\eqref{ext:eq-natural-conditional-solution}. Since $P\le1$, this
also gives $A(t)\ge\rho t$. Incorrect coordinates satisfy
$\dot p_i=-(c(P)P+\rho)p_i$ and remain positive at every finite time.
Correct conditional coordinates remain positive by the explicit solution,
and $P\ge P(0)>0$. Because $c$ is bounded on $[0,1]$ and
$\rho/P\le\rho/P(0)$, these formulas extend the interior solution to every
finite time.

For $\rho=0$, $P$ increases and cannot tend to a limit below one: at any
such limit $(1-P)c(P)P$ would stay positive. For $\rho>0$,
$\dot P\ge\rho(1-P)$ yields the correctness bound in
Equation~\eqref{ext:eq-natural-retention}. Each banked $q_b(t)$ lies between
$q_b(0)$ and $w_b$, so multiplication by $P(t)\ge P(0)$ proves the
probability floor. Finally, $A(t)\to\infty$ gives $q_c(t)\to\bar w_c$, so unbanked
correct modes have zero limiting mass.
\end{proof}

\paragraph{Worked example: the metric changes which modes grow.}
At the running example's $P(0)=0.5$, $q(0)=(0.6,0.3,0.1)$, and $G=4$,
Dr.GRPO has $c(P)=3/4$. Without replay, the Fisher flow gives
$\dot P(0)=0.1875$ and $\dot q(0)=0$: the model learns correctness while
keeping its correct-mode proportions. With the uniform three-mode target
$w=(1/3,1/3,1/3)$ and $\rho=0.1$, the same initial distribution instead gives
\[
 \dot P(0)=0.2375,\qquad
 \dot q(0)\approx(-0.05333,\ 0.00667,\ 0.04667).
\]
Replay now increases the rare third mode's share among correct responses.
Equation~\eqref{ext:eq-natural-retention} also guarantees $p_3(t)\ge0.05$
for all time in this ideal flow. These are instantaneous derivatives and a
trajectory bound, not the result of taking a discrete optimizer step.

The smooth potentials with positive derivatives in
Lemmas~\ref{lem:grpo-mean} and~\ref{lem:maxrl-mean} satisfy the premises. The Fisher-gradient formula
is established information geometry; the result above is its elementary
specialization to correctness selection plus replay toward a fixed target.
It makes the scope of the collapse claim concrete: increasing correctness
alone does not force a loss of conditional support under every optimization
geometry. The theorem concerns the exact categorical Fisher flow. It gives
no guarantee for approximate natural gradient, a shared neural Fisher
matrix, PPO, or AdamW, and its time variable is an ideal optimization clock
rather than elapsed training time.

